Using 36,528 aggregated 2023 Medicare DME records from the 50 U.S. states and Washington, DC, this study evaluated geographic prediction reliability using interpretable group-mean baselines. Distinguishing HCPCS codes by rental classification reduced actual-payment leave-one-area-out MAE from \$47.84 to \$19.89. Standardized average payments were also easier to predict under the tested baselines, with rental-aware MAE of \$10.30, and their area-level MAE was lower in all 51 held-out areas. Rare unseen HCPCS codes nonetheless accounted for a disproportionate share of absolute error under the HCPCS-only baseline. These findings demonstrate why payment-target definitions, equipment classification, code coverage, and area-specific errors should accompany pooled accuracy when evaluating geographic reliability. They do not establish causality or a universally superior prediction model; validation on additional years and modeling approaches is necessary.
